Nach dem dritten Kepler'schen Gesetz ist $a^3/T^2$ für alle Körper, die denselben Zentralkörper umkreisen, gleich (FoTaBe S.~160):
\begin{align}
 \frac{a_1^3}{T_1^2} &= \frac{a_2^3}{T_2^2}
 \quad\Rightarrow\quad \frac{a_1}{a_2} = \qaF
\end{align}
Mit $T_1/T_2 = \qTO$:
\begin{align}
 \frac{a_1}{a_2} &= \left(\qTO\right)^{2/3} = \qa \approx \result{\qaP{0}{2}}
\end{align}
Der Planet mit der doppelten Umlaufzeit hat eine rund \qaP{0}{2}-mal so grosse Halbachse.
